\section{CONCLUSIONS}

\bulletpoints{
    \item Topic Sentence: The take home message is that we were able to achieve real-time runtime while maintaining universality.
    \item Our approach is entirely numerical
    \item We achieved 120 microsecond runtime and maintained universality
    \item We also have a potential proxy
}

\commentblock{
This investigation covers explicit integration techniques, quaternion-based orientation representation, and efficient Jacobian updating schemes, resulting in significant reductions in the computational runtime of the forward kinematics of the Cosserat Rod model.
By combining the proposed optimizations, we demonstrate a reduction of up to $90\%$ in function calls with no loss in accuracy.
Our optimized kinematic model achieves an average runtime of just \SI{120}{\mu s}, making it suitable for high-frequency, real-time control of CTCR systems while also ensuring broad applicability across various architectures and systems.
Furthermore, we introduced a novel, computationally accessible proxy for elastic stability, which holds promise for detecting snapping behavior in CTCRs.
}

%\footnotetext{\url{https://github.com/ContinuumRoboticsLab/CRVisToolkit}}